\section{USAJMO 2019}
        \begin{enumerate}
        	\item (\textit{USAJMO 2019}) \textbf{Note:} For any geometry problem whose statement begins with an asterisk $(*)$, the first page of the solution must be a large, in-scale, clearly labeled diagram. Failure to meet this requirement will result in an automatic 1-point deduction.


 There are $a+b$ bowls arranged in a row, numbered $1$ through $a+b$, where $a$ and $b$ are given positive integers. Initially, each of the first $a$ bowls contains an apple, and each of the last $b$ bowls contains a pear.


 A legal move consists of moving an apple from bowl $i$ to bowl $i+1$ and a pear from bowl $j$ to bowl $j-1$, provided that the difference $i-j$ is even. We permit multiple fruits in the same bowl at the same time. The goal is to end up with the first $b$ bowls each containing a pear and the last $a$ bowls each containing an apple. Show that this is possible if and only if the product $ab$ is even.


 

	\item (\textit{USAJMO 2019}) \textbf{Note:} For any geometry problem whose statement begins with an asterisk $(*)$, the first page of the solution must be a large, in-scale, clearly labeled diagram. Failure to meet this requirement will result in an automatic 1-point deduction.


 Let $\mathbb Z$ be the set of all integers. Find all pairs of integers $(a,b)$ for which there exist functions $f:\mathbb Z\rightarrow\mathbb Z$ and $g:\mathbb Z\rightarrow\mathbb Z$ satisfying 
 \[f(g(x))=x+a\quad\text{and}\quad g(f(x))=x+b\] for all integers $x$.


 

	\item (\textit{USAJMO 2019}) \textbf{Note:} For any geometry problem whose statement begins with an asterisk $(*)$, the first page of the solution must be a large, in-scale, clearly labeled diagram. Failure to meet this requirement will result in an automatic 1-point deduction.


 $(*)$  Let $ABCD$ be a cyclic quadrilateral satisfying $AD^2+BC^2=AB^2$. The diagonals of $ABCD$ intersect at $E$. Let $P$ be a point on side $\overline{AB}$ satisfying $\angle APD=\angle BPC$. Show that line $PE$ bisects $\overline{CD}$.


 



 

	\item (\textit{USAJMO 2019}) \textbf{Note:} For any geometry problem whose statement begins with an asterisk $(*)$, the first page of the solution must be a large, in-scale, clearly labeled diagram. Failure to meet this requirement will result in an automatic 1-point deduction.


 $(*)$ Let $ABC$ be a triangle with $\angle ABC$ obtuse. The $A$\textit{-excircle} is a circle in the exterior of $\triangle ABC$ that is tangent to side $BC$ of the triangle and tangent to the extensions of the other two sides. Let $E, F$ be the feet of the altitudes from $B$ and $C$ to lines $AC$ and $AB$, respectively. Can line $EF$ be tangent to the $A$-excircle?


 

	\item (\textit{USAJMO 2019}) \textbf{Note:} For any geometry problem whose statement begins with an asterisk $(*)$, the first page of the solution must be a large, in-scale, clearly labeled diagram. Failure to meet this requirement will result in an automatic 1-point deduction.


 Let $n$ be a nonnegative integer. Determine the number of ways that one can choose $(n+1)^2$ sets $S_{i,j} \subseteq \{1,2,...,2n\}$, for integers $i,j$ with $0 \le i,j \le n$ such that:


 $\bullet$ for all $0 \le i,j \le n$, the set $S_{i,j}$ has $i+j$ elements; and


 $\bullet$ $S_{i,j} \subseteq S_{k,l}$ whenever $0 \le i \le k \le n$ and $0 \le j \le l \le n$


 

	\item (\textit{USAJMO 2019}) \textbf{Note:} For any geometry problem whose statement begins with an asterisk $(*)$, the first page of the solution must be a large, in-scale, clearly labeled diagram. Failure to meet this requirement will result in an automatic 1-point deduction.


 Two rational numbers $\frac{m}{n}$ and $\frac{n}{m}$ are written on a blackboard, where $m$ and $n$ are relatively prime positive integers. At any point, Evan may pick two of the numbers $x$ and $y$ written on the board and write either their arithmetic mean $\frac{x+y}{2}$ or their harmonic mean $\frac{2xy}{x+y}$ on the board as well. Find all pairs $(m,n)$ such that Evan can write $1$ on the board in finitely many steps.


 

\end{enumerate}